% GENERATED FILE. Edit canonical lessons, then run npm run generate:books.
% canonical-source-hash: fnv1a64:3f1715d617eb5fce
% canonical-lessons: TA-C556-cirukatai, TA-C556-ilakkiyam, TA-C556-kaivinai, TA-C556-necavu, TA-C556-carikai

\chapter{A Writing and Craft Contest}
\label{ch:ta-a-writing-and-craft-contest}
\hlchaptermodality{556}

\begin{chapteropening}
\textbf{By the end of this chapter:} \emph{I can talk about the rules of a contest.}

Complete the last lesson of chapter 556: I can talk about the rules of a contest.
\end{chapteropening}

\section[\texorpdfstring{\ta{சிறுகதை} (ciṟukatai)}{சிறுகதை (ciṟukatai)}]{\ta{சிறுகதை} (ciṟukatai) — a short story}

\label{lesson:TA-C556-cirukatai}



\begin{quote}
Before the new one: say the Tamil for a review, then the Tamil for a story character.
\end{quote}

\subsection*{You'll want to know: \ta{சிறுகதை}}
\textbf{\ta{சிறுகதை}} — \emph{ciṟukatai} — \textquotedblleft{}a short story\textquotedblright{}.

\begin{quote}
\ta{உங்கள்} \ta{சிறுகதை} \ta{தேர்வானால்} \ta{இதழில்} \ta{வெளிவரும்}.
\end{quote}

\begin{quote}
\emph{uṅkaḷ ciṟukatai tērvāṉāl itaḻil veḷivarum.}
\end{quote}

\begin{quote}
If your short story is chosen, it will appear in the magazine.
\end{quote}

\ta{சிறு} \textquotedblleft{}small\textquotedblright{} + \ta{கதை} (\emph{katai}). The plural is \ta{சிறுகதைகள்}.

\begin{cousinweb}
\textbf{\ta{சிறுகதை}} (\emph{ciṟukatai}) joins the native \ta{சிறு} (\emph{ciṟu}), \textquotedblleft{}small\textquotedblright{} (DEDR 1594), and \ta{கதை} (\emph{katai}), \textquotedblleft{}a story\textquotedblright{}, from Sanskrit \emph{kathā} (CDIAL 2704).
\end{cousinweb}

\begin{grammarlens}[title={What you've built}]
One more word for saying what happens if.
\end{grammarlens}

\subsection*{Your turn}
\begin{itemize}
\raggedright
  \item \emph{Say it:} \emph{ciṟukatai}
  \item \emph{Say it:} \emph{ciṟukatai}, once more
  \item \emph{Say it:} \emph{uṅkaḷ ciṟukatai tērvāṉāl itaḻil veḷivarum.} — If your short story is chosen, it will appear in the magazine.
  \item \emph{Recall:} say the Tamil for a review, then the Tamil for a story character, then say \emph{ciṟukatai} again
\end{itemize}

\begin{tcolorbox}[breakable,colback=teal!4,colframe=teal!35!black,title={Before you move on}]
What does \ta{சிறுகதை} mean? (\textquotedblleft{}A short story\textquotedblright{}.) Say it once more.
\end{tcolorbox}
\section[\texorpdfstring{\ta{இலக்கியம்} (ilakkiyam)}{இலக்கியம் (ilakkiyam)}]{\ta{இலக்கியம்} (ilakkiyam) — literature}

\label{lesson:TA-C556-ilakkiyam}



\begin{quote}
Before the new one: say the Tamil for a story character, then the Tamil for a short story.
\end{quote}

\subsection*{You'll want to know: \ta{இலக்கியம்}}
\textbf{\ta{இலக்கியம்}} — \emph{ilakkiyam} — \textquotedblleft{}literature\textquotedblright{}.

\begin{quote}
\ta{தமிழ்} \ta{இலக்கியம்} \ta{இரண்டாயிரம்} \ta{ஆண்டுகளுக்கும்} \ta{மேல்} \ta{பழமையானது}.
\end{quote}

\begin{quote}
\emph{tamiḻ ilakkiyam iraṇṭāyiram āṇṭukaḷukkum mēl paḻamaiyāṉatu.}
\end{quote}

\begin{quote}
Tamil literature is more than two thousand years old.
\end{quote}

A -\ta{அம்} noun: \ta{இலக்கியத்தில்}. An \ta{இலக்கியவாதி} is a literary man.

\begin{cousinweb}
\textbf{\ta{இலக்கியம்}} (\emph{ilakkiyam}) is usually traced to Sanskrit \emph{lakṣya}, \textquotedblleft{}a mark, something kept in view\textquotedblright{} (Monier-Williams), the source of \ta{லட்சியம்} (\emph{laṭciyam}) too; Tamil puts \ta{இ} before the initial \ta{ல}. How it came to mean literature is not settled.
\end{cousinweb}

\begin{grammarlens}[title={What you've built}]
One more word for saying what happens if.
\end{grammarlens}

\subsection*{Your turn}
\begin{itemize}
\raggedright
  \item \emph{Say it:} \emph{ilakkiyam}
  \item \emph{Say it:} \emph{ilakkiyam}, once more
  \item \emph{Say it:} \emph{tamiḻ ilakkiyam iraṇṭāyiram āṇṭukaḷukkum mēl paḻamaiyāṉatu.} — Tamil literature is more than two thousand years old.
  \item \emph{Recall:} say the Tamil for a story character, then the Tamil for a short story, then say \emph{ilakkiyam} again
\end{itemize}

\begin{tcolorbox}[breakable,colback=teal!4,colframe=teal!35!black,title={Before you move on}]
What does \ta{இலக்கியம்} mean? (\textquotedblleft{}Literature\textquotedblright{}.) Say it once more.
\end{tcolorbox}
\section[\texorpdfstring{\ta{கைவினை} (kaiviṉai)}{கைவினை (kaiviṉai)}]{\ta{கைவினை} (kaiviṉai) — a craft; handicraft}

\label{lesson:TA-C556-kaivinai}



\begin{quote}
Before the new one: say the Tamil for a short story, then the Tamil for literature.
\end{quote}

\subsection*{You'll want to know: \ta{கைவினை}}
\textbf{\ta{கைவினை}} — \emph{kaiviṉai} — \textquotedblleft{}a craft; handicraft\textquotedblright{}.

\begin{quote}
\ta{கைவினைப்} \ta{பொருட்களைப்} \ta{போட்டிக்கு} \ta{அனுப்பினால்} \ta{பரிசு} \ta{கிடைக்கலாம்}.
\end{quote}

\begin{quote}
\emph{kaiviṉaip poruṭkaḷaip pōṭṭikku aṉuppiṉāl paricu kiṭaikkalām.}
\end{quote}

\begin{quote}
If you send craftwork to the contest, you may win a prize.
\end{quote}

Craftwork is \ta{கைவினைப்} \ta{பொருட்கள்}: \ta{கை} \textquotedblleft{}hand\textquotedblright{} + \ta{வினை} \textquotedblleft{}work\textquotedblright{}.

\begin{cousinweb}
\textbf{\ta{கைவினை}} (\emph{kaiviṉai}) is a native compound: \ta{கை} (\emph{kai}), \textquotedblleft{}a hand\textquotedblright{}, and \ta{வினை} (\emph{viṉai}), \textquotedblleft{}action, work\textquotedblright{} (DEDR 5445).
\end{cousinweb}

\begin{grammarlens}[title={What you've built}]
One more word for saying what happens if.
\end{grammarlens}

\subsection*{Your turn}
\begin{itemize}
\raggedright
  \item \emph{Say it:} \emph{kaiviṉai}
  \item \emph{Say it:} \emph{kaiviṉai}, once more
  \item \emph{Say it:} \emph{kaiviṉaip poruṭkaḷaip pōṭṭikku aṉuppiṉāl paricu kiṭaikkalām.} — If you send craftwork to the contest, you may win a prize.
  \item \emph{Recall:} say the Tamil for a short story, then the Tamil for literature, then say \emph{kaiviṉai} again
\end{itemize}

\begin{tcolorbox}[breakable,colback=teal!4,colframe=teal!35!black,title={Before you move on}]
What does \ta{கைவினை} mean? (\textquotedblleft{}A craft; handicraft\textquotedblright{}.) Say it once more.
\end{tcolorbox}
\section[\texorpdfstring{\ta{நெசவு} (necavu)}{நெசவு (necavu)}]{\ta{நெசவு} (necavu) — weaving}

\label{lesson:TA-C556-necavu}



\begin{quote}
Before the new one: say the Tamil for literature, then the Tamil for a craft.
\end{quote}

\subsection*{You'll want to know: \ta{நெசவு}}
\textbf{\ta{நெசவு}} — \emph{necavu} — \textquotedblleft{}weaving\textquotedblright{}.

\begin{quote}
\ta{நெசவு} \ta{எங்கள்} \ta{கிராமத்தின்} \ta{முக்கியத்} \ta{தொழில்}.
\end{quote}

\begin{quote}
\emph{necavu eṅkaḷ kirāmattiṉ mukkiyat toḻil.}
\end{quote}

\begin{quote}
Weaving is the main trade of our village.
\end{quote}

The verb is \ta{நெய்} (\emph{ney}), to weave; a weaver is a \ta{நெசவாளர்} (\emph{necavāḷar}).

\begin{cousinweb}
\textbf{\ta{நெசவு}} (\emph{necavu}) is native Dravidian: DEDR 3745 gives \ta{நெய்} (\emph{ney}), \textquotedblleft{}to weave\textquotedblright{}, and \ta{நெசவு} (\emph{necavu}), \textquotedblleft{}weaving, a web\textquotedblright{}.
\end{cousinweb}

\begin{grammarlens}[title={What you've built}]
One more word for saying what happens if.
\end{grammarlens}

\subsection*{Your turn}
\begin{itemize}
\raggedright
  \item \emph{Say it:} \emph{necavu}
  \item \emph{Say it:} \emph{necavu}, once more
  \item \emph{Say it:} \emph{necavu eṅkaḷ kirāmattiṉ mukkiyat toḻil.} — Weaving is the main trade of our village.
  \item \emph{Recall:} say the Tamil for literature, then the Tamil for a craft, then say \emph{necavu} again
\end{itemize}

\begin{tcolorbox}[breakable,colback=teal!4,colframe=teal!35!black,title={Before you move on}]
What does \ta{நெசவு} mean? (\textquotedblleft{}Weaving\textquotedblright{}.) Say it once more.
\end{tcolorbox}
\section[\texorpdfstring{\ta{சரிகை} (carikai)}{சரிகை (carikai)}]{\ta{சரிகை} (carikai) — gold thread; zari}

\label{lesson:TA-C556-carikai}



\begin{quote}
Before the new one: say the Tamil for a craft, then the Tamil for weaving.
\end{quote}

\subsection*{You'll want to know: \ta{சரிகை}}
\textbf{\ta{சரிகை}} — \emph{carikai} — \textquotedblleft{}gold thread; zari\textquotedblright{}.

\begin{quote}
\ta{சேலையின்} \ta{ஓரத்தில்} \ta{சரிகை} \ta{இருந்தால்} \ta{விலை} \ta{அதிகம்}.
\end{quote}

\begin{quote}
\emph{cēlaiyiṉ ōrattil carikai iruntāl vilai atikam.}
\end{quote}

\begin{quote}
If there is gold thread on the border of the sari, the price is higher.
\end{quote}

The gold or silver thread woven into the border of a silk sari.

\begin{cousinweb}
The origin of \textbf{\ta{சரிகை}} (\emph{carikai}) is not settled: it may be the same word as \emph{zarī} in other Indian languages, but no source we can check shows the path.
\end{cousinweb}

\begin{grammarlens}[title={What you've built}]
One more word for saying what happens if.
\end{grammarlens}

\subsection*{Your turn}
\begin{itemize}
\raggedright
  \item \emph{Say it:} \emph{carikai}
  \item \emph{Say it:} \emph{carikai}, once more
  \item \emph{Say it:} \emph{cēlaiyiṉ ōrattil carikai iruntāl vilai atikam.} — If there is gold thread on the border of the sari, the price is higher.
  \item \emph{Recall:} say the Tamil for a craft, then the Tamil for weaving, then say \emph{carikai} again
\end{itemize}

\begin{tcolorbox}[breakable,colback=teal!4,colframe=teal!35!black,title={Before you move on}]
What does \ta{சரிகை} mean? (\textquotedblleft{}Gold thread; zari\textquotedblright{}.) Say it once more.
\end{tcolorbox}
